\chapter{Developer Documentation}
\label{ch:developer}

This chapter provides comprehensive technical documentation for developers working with the Solar Energy Monitoring and Control System (Nebula Power). It covers the solution plan including system architecture and design, implementation details with development decisions and component structure, comprehensive testing strategies including scalability analysis, and development environment setup. The documentation is designed to help developers understand, extend, and maintain the system effectively while adhering to academic standards for technical documentation.

This documentation is structured into three primary sections: \textbf{Solution Plan} (Section~\ref{sec:solution-plan}) which details the architectural design and planning decisions, \textbf{Implementation} (Section~\ref{sec:implementation}) which covers the technical implementation and development approach, and \textbf{Testing} (Section~\ref{sec:testing-comprehensive}) which provides comprehensive testing strategies and analysis.

\section{Solution Plan}
\label{sec:solution-plan}

This section presents the comprehensive solution plan for the Solar Energy Monitoring and Control System, including system architecture design, database design, modular structure planning, and user interface design. The solution plan follows academic standards for technical system design documentation and provides the foundational framework for the entire system implementation.

\subsection{System Architecture Design}
\label{subsec:system-architecture-design}

The Solar Energy Monitoring and Control System employs a sophisticated multi-layered architecture that balances modern software engineering principles with practical deployment considerations. The architecture follows a modular monolithic pattern \cite{martin2019clean}, providing clear separation of concerns while maintaining deployment simplicity and operational efficiency \cite{newman2021building}.

\subsubsection{Architectural Overview and Design Rationale}
\label{subsubsec:architecture-overview-rationale}

The system architecture is designed around five core principles: \textbf{modularity} for maintainability, \textbf{scalability} for growth accommodation, \textbf{security} for data protection, \textbf{reliability} for continuous operation, and \textbf{performance} for real-time processing requirements.

\begin{figure}[h]
    \centering
    \tiny
    \begin{verbatim}
+---------------------------------------------------------------+
|                      CLIENT APPLICATIONS                      |
+---------------------------------------------------------------+
|  Web Frontend (Next.js/React)                                 |
|  - Admin Dashboard, Customer Portal                           |
|  - Real-time Monitoring via STOMP/SockJS                      |
+---------------------------------------------------------------+
                              |
                          HTTPS / WSS                           |
                              |
+---------------------------------------------------------------+
|                       APPLICATION LAYER                       |
+---------------------------------------------------------------+
| User Mgmt     | Energy Monitoring | Payment Compliance        |
| Auth/Profile  | Readings, Series  | Plans, Payments, Reports  |
|               | Summaries, WS     | Reminders, Compliance     |
+---------------------------------------------------------------+
| Tamper & Security         | Service Control & Integration    |
| Events, Logs, Detection   | Status, Commands, Heartbeats     |
+---------------------------------------------------------------+
                              |
                           REST / WS                            |
                              |
+---------------------------------------------------------------+
|                        PERSISTENCE LAYER                      |
+---------------------------------------------------------------+
| PostgreSQL / H2 (dev)                                         |
| Entities: users, solar_installations, energy_data, summaries, |
| payments, payment_plans, reminders, service_status, commands, |
| operational_logs, tamper_events, security_logs, configs       |
+---------------------------------------------------------------+
    \end{verbatim}
    \caption{Comprehensive System Architecture Overview}
    \label{fig:system-architecture-comprehensive}
\end{figure}

The system architecture is illustrated in Figure~\ref{fig:system-architecture-comprehensive}, which provides a comprehensive overview of all layers and their interactions.

The system uses a modular monolithic architecture with clear separation between frontend, backend, and database. Communication is via REST APIs and WebSockets. Core modules are User Management, Energy Monitoring, Payment Compliance, Security, and Service Control.

\paragraph{Technology Stack Selection and Rationale}

The technology stack selection was driven by specific architectural requirements and long-term maintainability considerations, as outlined in Table~\ref{tab:tech-stack-rationale}:

\begin{table}[h]
    \centering
    \caption{Technology Stack Selection with Rationale}
    \label{tab:tech-stack-rationale}
    \begin{tabular}{|p{3.2cm}|p{4cm}|p{6.3cm}|}
        \hline
        \textbf{Layer} & \textbf{Technology} & \textbf{Selection Rationale} \\
        \hline
        Backend Core & Java 21, Spring Boot 3.3.4 & Type safety, mature ecosystem, excellent security framework integration, performance characteristics suitable for real-time processing \\
        \hline
        Frontend & Next.js 15, React 19, TypeScript & Server-side rendering for performance, component reusability, type safety, excellent developer experience \\
        \hline
        Database & PostgreSQL 15, H2 (testing) & ACID compliance, JSON support, excellent performance with time-series data, robust backup/recovery \\
        \hline
        Security & Spring Security, JWT, BCrypt & Industry-standard authentication, stateless tokens, proven encryption algorithms \\
        \hline
        Communication & REST APIs, WebSocket, HTTP/2 & Standardized protocols, real-time capabilities, efficient data transfer \\
        \hline
        Containerization & Docker, Docker Compose & Consistent deployment environments, scalability, development/production parity \\
        \hline
        Testing & JUnit 5, Mockito, Jest, React Testing Library & Comprehensive testing capabilities, mocking support, component testing \\
        \hline
    \end{tabular}
\end{table}

\subsubsection{Modular Architecture Design}
\label{subsubsec:modular-architecture-design}

The system implements a carefully planned modular architecture with five core modules, each designed with specific responsibilities and clear interfaces, as shown in Figures~\ref{fig:module-architecture-part1} and~\ref{fig:module-architecture-part2}:

\begin{figure}[h]
    \centering
    \begin{minipage}[t]{0.48\textwidth}
        \centering
        \small
        \begin{verbatim}
USER MANAGEMENT MODULE
+------------------------------+
| Controllers/Services         |
| +-- AuthController           |
| +-- UserProfileController    |
| +-- CustomerController       |
| +-- SecurityService          |
| +-- UserService              |
| +-- UserRepository           |
| +-- SecurityConfig           |
|                              |
| Responsibilities:            |
| +-- JWT authentication       |
| +-- User profile mgmt        |
| +-- Password security        |
+------------------------------+

ENERGY MONITORING MODULE
+------------------------------+
| Controllers/Services         |
| +-- EnergyDataController     |
| +-- EnergySummaryController  |
| +-- SolarInstallationController |
| +-- WebSocketService         |
|                              |
| Responsibilities:            |
| +-- Real-time monitoring     |
| +-- Installation management  |
| +-- Efficiency calculations  |
| +-- Historical aggregation   |
+------------------------------+

PAYMENT COMPLIANCE MODULE
+------------------------------+
| Controllers/Services         |
| +-- CustomerPaymentController|
| +-- AdminPaymentController   |
| +-- PaymentReportController  |
| +-- Payment/Plan/ReminderSvc |
|                              |
| Responsibilities:            |
| +-- Payment plan creation    |
| +-- Payment processing       |
| +-- Automated billing        |
+------------------------------+
        \end{verbatim}
        \caption{Core Business Modules}
        \label{fig:module-architecture-part1}
    \end{minipage}
    \hfill
    \begin{minipage}[t]{0.48\textwidth}
        \centering
        \small
        \begin{verbatim}
TAMPERING DETECTION MODULE
+------------------------------+
| Controllers/Services         |
| +-- TamperEventController    |
| +-- TamperDetectionController|
| +-- SecurityLogController    |
|                              |
| Responsibilities:            |
| +-- Tamper event detection   |
| +-- Security alerts          |
| +-- Incident response        |
| +-- Security analytics       |
| +-- Compliance monitoring    |
+------------------------------+

SERVICE CONTROL MODULE
+------------------------------+
| Controllers/Services         |
| +-- ServiceStatusController  |
| +-- DeviceCommandController  |
| +-- SystemIntegrationController |
| +-- IntegrationController    |
|                              |
|                              |
| Responsibilities:            |
| +-- Device health monitoring |
| +-- Command execution        |
| +-- System diagnostics       |
| +-- Operational logging      |
| +-- Config management        |
+------------------------------+
        \end{verbatim}
        \caption{Security and Control Modules}
        \label{fig:module-architecture-part2}
    \end{minipage}
\end{figure}

\paragraph{Short Code Excerpts (verified)}
\begin{verbatim}
// EnergyDataController (extract)
@PostMapping("/readings")
public ResponseEntity<EnergyDataDTO> submitEnergyReading(
    @Valid @RequestBody EnergyDataRequest request) {
  return ResponseEntity.ok(energyDataService.processEnergyData(request));
}

// WebSocketConfig (extract)
config.enableSimpleBroker("/topic");
registry.addEndpoint("/ws").setAllowedOriginPatterns("*").withSockJS();

// SecurityService (extract)
public boolean hasAccessToInstallation(Long installationId) {
  Authentication authentication = SecurityContextHolder.getContext().getAuthentication();
  if (authentication == null || !authentication.isAuthenticated()) return false;
  if (authentication.getAuthorities().stream()
        .anyMatch(a -> a.getAuthority().equals("ROLE_ADMIN"))) return true;
  Object principal = authentication.getPrincipal();
  if (principal instanceof UserPrincipal) {
    Long userId = ((UserPrincipal) principal).getId();
    return installationRepository.findById(installationId)
      .map(i -> i.getUser() != null && i.getUser().getId().equals(userId))
      .orElse(false);
  }
  return false;
}
\end{verbatim}

\subsubsection{Inter-Module Communication Design}
\label{subsubsec:inter-module-communication}

The system implements a sophisticated communication strategy that balances performance with maintainability:

\paragraph{Communication Patterns}
\begin{itemize}
    \item \textbf{Direct Service Calls}: For synchronous operations requiring immediate response
    \item \textbf{Spring Events}: For asynchronous, loosely-coupled communication within the application context
    \item \textbf{Shared Data Access}: Through standardized repository interfaces for data consistency
    \item \textbf{REST API Interfaces}: For potential future microservice decomposition
\end{itemize}

% --- Replace detailed communication flow with concise summary ---

Inter-module communication is handled via direct service calls for synchronous operations and Spring Events for asynchronous updates. Data consistency is maintained through shared repository interfaces. REST APIs are used for frontend-backend communication.

Example: When a device submits an energy reading, the backend processes the data, updates analytics, and triggers compliance checks if needed.

% --- End concise summary ---

\subsection{Database Design and Architecture}
\label{subsec:database-design}

The database design reflects the implemented JPA model and focuses on transactional integrity and straightforward time‑series access for energy data.

\subsubsection{Database Schema Design}
\label{subsubsec:database-schema}

Domains and key entities (tables):
\begin{itemize}
  \item \textbf{User}: \texttt{users} (id,email,password,role,enabled,...)
  \item \textbf{Energy}: \texttt{solar\_installations}, \texttt{energy\_data}, \texttt{energy\_summaries}
  \item \textbf{Payments}: \texttt{payment\_plans}, \texttt{payments}, \texttt{payment\_reminders}, \texttt{reminder\_configs}, \texttt{grace\_period\_configs}
  \item \textbf{Security}: \texttt{tamper\_events}, \texttt{security\_logs}, \texttt{tamper\_monitoring\_status}, \texttt{alert\_configs}
  \item \textbf{Service Control}: \texttt{service\_status}, \texttt{device\_commands}, \texttt{operational\_logs}
\end{itemize}

Relationships (selected): \texttt{energy\_data}$\rightarrow$\texttt{solar\_installations} (many-to-one), \texttt{payments}$\rightarrow$\texttt{solar\_installations} (many-to-one), \texttt{payments}$\rightarrow$\texttt{payment\_plans} (many-to-one), \texttt{tamper\_events}$\rightarrow$\texttt{solar\_installations} (many-to-one), \texttt{alert\_configs}$\leftrightarrow$\texttt{solar\_installations} (one-to-one).

Indexing: a composite index exists on \texttt{energy\_data(installation\_id,timestamp)} to accelerate time‑range queries per installation. Consider additional indexes on foreign keys and date/status columns in high‑traffic tables (e.g., \texttt{payments}, \texttt{service\_status}).

\subsubsection{Database Performance Optimization}
\label{subsubsec:database-optimization}

The database design incorporates several performance optimization strategies:

\paragraph{Indexing Strategy}
\begin{itemize}
  \item Composite index on \texttt{energy\_data(installation\_id,timestamp)} (implemented)
  \item Recommended: indexes on \texttt{payments(installation\_id,due\_date)} and \texttt{service\_status(installation\_id)}
\end{itemize}

\paragraph{Partitioning}
Not configured in this implementation; may be considered for very large datasets.

\paragraph{Data Retention Policy}
\begin{itemize}
    \item \textbf{Hot Data}: Last 30 days in primary tables with full indexing
    \item \textbf{Warm Data}: 31-365 days in compressed partitions
    \item \textbf{Cold Data}: > 1 year archived to separate storage with minimal indexing
\end{itemize}

\subsection{User Interface Design and Architecture}
\label{subsec:ui-design-architecture}

The user interface design follows modern UX/UI principles with a focus on role-based functionality, responsive design, and real-time data visualization. The frontend architecture employs Next.js 15 with React 19 and TypeScript for type safety and maintainability.

\subsubsection{Frontend Architecture Design}
\label{subsubsec:frontend-architecture}

The frontend architecture implements a component-based design with clear separation of concerns, as illustrated in Figures~\ref{fig:frontend-architecture-part1} and \ref{fig:frontend-architecture-part2}:

\begin{figure}[htbp]
    \centering
    \begin{minipage}[t]{0.48\textwidth}
        \centering
        \scriptsize
        \begin{verbatim}
PRESENTATION & COMPONENT LAYERS

+------------------------------+
|     PRESENTATION LAYER       |
+------------------------------+
| Pages (Next.js App Router)  |
| +-- app/                     |
| |   +-- (auth)/              |
| |   |   +-- login/page.tsx   |
| |   |   +-- register/page.tsx|
| |   +-- admin/               |
| |   |   +-- dashboard/       |
| |   |   +-- customers/       |
| |   |   +-- monitoring/      |
| |   |   +-- payments/        |
| |   |   +-- security/        |
| |   |   +-- layout.tsx       |
| |   +-- customer/            |
| |   |   +-- dashboard/       |
| |   |   +-- energy/          |
| |   |   +-- payments/        |
| |   |   +-- system/          |
| |   |   +-- layout.tsx       |
| |   +-- layout.tsx           |
| |   +-- globals.css          |
| |   +-- page.tsx             |
+------------------------------+

+------------------------------+
|      COMPONENT LAYER         |
+------------------------------+
| UI Components (components/)  |
| +-- Layout Components        |
| |   +-- Header.tsx           |
| |   +-- Sidebar.tsx          |
| |   +-- Footer.tsx           |
| |   +-- Navigation.tsx       |
| +-- Form Components          |
| |   +-- Input.tsx            |
| |   +-- Button.tsx           |
| |   +-- Select.tsx           |
| |   +-- DatePicker.tsx       |
| |   +-- FileUpload.tsx       |
| +-- Data Display             |
| |   +-- DataTable.tsx        |
| |   +-- Card.tsx             |
| |   +-- Badge.tsx            |
| |   +-- Progress.tsx         |
| |   +-- Tooltip.tsx          |
| +-- Charts & Visualizations  |
| |   +-- LineChart.tsx        |
| |   +-- BarChart.tsx         |
| |   +-- PieChart.tsx         |
| |   +-- GaugeChart.tsx       |
| |   +-- SparklineChart.tsx   |
| +-- Feedback Components      |
|     +-- Modal.tsx            |
|     +-- Toast.tsx            |
|     +-- LoadingSpinner.tsx   |
|     +-- ErrorBoundary.tsx    |
|     +-- Alert.tsx            |
+------------------------------+
        \end{verbatim}
        \caption{Frontend UI and Presentation Layers}
        \label{fig:frontend-architecture-part1}
    \end{minipage}
    \hfill
    \begin{minipage}[t]{0.48\textwidth}
        \centering
        \scriptsize
        \begin{verbatim}
STATE & SERVICE LAYERS

+------------------------------+
|       STATE LAYER            |
+------------------------------+
| State Management             |
| +-- Contexts                 |
| |   +-- AuthContext.tsx      |
| |   +-- ThemeContext.tsx     |
| |   +-- NotificationContext  |
| |   +-- WebSocketContext.tsx |
| +-- Custom Hooks             |
| |   +-- useAuth.ts           |
| |   +-- useApi.ts            |
| |   +-- useWebSocket.ts      |
| |   +-- useLocalStorage.ts   |
| |   +-- useDebounce.ts       |
| +-- Utilities                |
|     +-- api.ts               |
|     +-- constants.ts         |
|     +-- types.ts             |
|     +-- utils.ts             |
|     +-- validation.ts        |
+------------------------------+

+------------------------------+
|      SERVICE LAYER           |
+------------------------------+
| API Services (lib/services/) |
| +-- authService.ts           |
| +-- energyService.ts         |
| +-- paymentService.ts        |
| +-- securityService.ts       |
| +-- systemService.ts         |
| +-- userService.ts           |
| +-- websocketService.ts      |
+------------------------------+

COMMUNICATION FLOW:
+-------------+
| Pages       |
+-----+-------+
      |
+-----v-------+
| Components  |
+-----+-------+
      |
+-----v-------+
| State/Hooks |
+-----+-------+
      |
+-----v-------+
| API Services|
+-------------+
        \end{verbatim}
        \caption{Frontend State and Service Layers}
        \label{fig:frontend-architecture-part2}
    \end{minipage}
\end{figure}

\subsubsection{Role-Based User Experience Design}
\label{subsubsec:role-based-ux}

The system implements distinct user experiences for different user roles:

\paragraph{Customer Interface Design}
\begin{itemize}
    \item \textbf{Dashboard Focus}: Real-time energy generation and consumption metrics
    \item \textbf{Simplified Navigation}: Task-oriented navigation with clear visual hierarchy
    \item \textbf{Energy Visualization}: Interactive charts showing production trends and efficiency
    \item \textbf{Payment Management}: Streamlined payment processing and history tracking
    \item \textbf{Alert Center}: Non-intrusive notifications for system status and payments
\end{itemize}

\paragraph{Administrator Interface Design}
\begin{itemize}
    \item \textbf{Comprehensive Dashboard}: System-wide metrics and health indicators
    \item \textbf{Advanced Navigation}: Multi-level navigation with quick access to all functions
    \item \textbf{Management Tools}: Customer management, system configuration, and reporting
    \item \textbf{Analytics Interface}: Advanced data visualization and reporting capabilities
    \item \textbf{Control Center}: System administration and security management tools
\end{itemize}

\subsubsection{Responsive Design Strategy}
\label{subsubsec:responsive-design-strategy}

The responsive design implementation follows a mobile-first approach with progressive enhancement, as detailed in Table~\ref{tab:responsive-breakpoints}:

\begin{table}[h]
    \centering
    \caption{Responsive Design Breakpoint Strategy}
    \label{tab:responsive-breakpoints}
    \begin{tabular}{|p{2cm}|p{2.5cm}|p{8.5cm}|}
        \hline
        \textbf{Breakpoint} & \textbf{Target Device} & \textbf{Design Adaptations} \\
        \hline
        < 640px & Mobile Phones & Single-column layout, collapsible navigation, touch-optimized controls, simplified charts \\
        \hline
        640px - 768px & Large Phones, Small Tablets & Two-column layout, side navigation drawer, enhanced touch targets \\
        \hline
        768px - 1024px & Tablets, Small Laptops & Multi-column layout, persistent sidebar, full chart functionality \\
        \hline
        1024px - 1280px & Laptops, Desktops & Full desktop layout, advanced features, keyboard navigation \\
        \hline
        > 1280px & Large Displays & Wide layouts, multi-panel views, enhanced data density \\
        \hline
    \end{tabular}
\end{table}

\subsubsection{Real-Time Data Visualization}
\label{subsubsec:realtime-visualization}

The system implements sophisticated real-time data visualization using WebSocket connections and optimized rendering:

\begin{itemize}
    \item \textbf{Live Energy Metrics}: Real-time power generation and consumption displays
    \item \textbf{Dynamic Charts}: Auto-updating time-series charts with configurable time ranges
    \item \textbf{Status Indicators}: Live device status and health monitoring displays
    \item \textbf{Alert Notifications}: Real-time security and system alerts with visual indicators
    \item \textbf{Performance Optimization}: Efficient rendering with data throttling and virtualization
\end{itemize}

\subsection{Security Architecture and Design}
\label{subsec:security-architecture-design}

The security architecture implements a comprehensive defense-in-depth strategy with multiple layers of protection, following industry best practices and compliance standards including GDPR, CCPA, and OWASP Top 10 recommendations.

\subsubsection{Multi-Layered Security Design}
\label{subsubsec:multi-layered-security}

The comprehensive security architecture is shown in Figures~\ref{fig:security-architecture-part1} and~\ref{fig:security-architecture-part2}, which illustrate the multiple layers of protection:

\begin{figure}[htbp]
    \centering
    \begin{minipage}[t]{0.48\textwidth}
        \centering
        \scriptsize
        \begin{verbatim}
NETWORK & APPLICATION SECURITY

+-------------------------+
| NETWORK SECURITY LAYER  |
+-------------------------+
| +-- TLS 1.3 Encryption  |
| +-- CORS Policy Config  |
| +-- Rate Limiting       |
|     (100 req/min/IP)    |
| +-- DDoS Protection     |
| +-- Firewall Rules      |
| +-- Network Segmentation|
+-------------------------+

+-------------------------+
| APPLICATION SECURITY    |
+-------------------------+
| Authentication & Auth:  |
| +-- JWT with RS256      |
| +-- Role-based access   |
|     (ADMIN, CUSTOMER)   |
| +-- Method-level        |
|     security            |
| +-- URL pattern         |
|     restrictions        |
| +-- Session management  |
|                         |
| Input Validation:       |
| +-- Bean Validation     |
|     (JSR-303)          |
| +-- SQL Injection       |
|     prevention         |
| +-- XSS Protection      |
| +-- CSRF Protection     |
| +-- File upload         |
|     validation         |
|                         |
| API Security:           |
| +-- OpenAPI 3.0 spec    |
| +-- Request/Response    |
|     validation         |
| +-- Security headers    |
|     (HSTS, CSP)        |
| +-- API versioning      |
| +-- Audit logging       |
+-------------------------+
        \end{verbatim}
        \caption{Network and Application Security Layers}
        \label{fig:security-architecture-part1}
    \end{minipage}
    \hfill
    \begin{minipage}[t]{0.48\textwidth}
        \centering
        \scriptsize
        \begin{verbatim}
DATA, DEVICE & MONITORING

+-------------------------+
| DATA SECURITY LAYER     |
+-------------------------+
| Encryption:             |
| +-- AES-256 at rest     |
| +-- BCrypt passwords    |
|     (12 salt rounds)    |
| +-- DB connection       |
|     encryption         |
| +-- Key management      |
|                         |
| Data Privacy:           |
| +-- GDPR compliance     |
| +-- CCPA compliance     |
| +-- Data anonymization  |
| +-- Right to deletion   |
| +-- Data export         |
|                         |
| Database Security:      |
| +-- User role           |
|     segregation        |
| +-- Connection pooling  |
| +-- Audit logging       |
| +-- Backup encryption   |
+-------------------------+

+-------------------------+
| DEVICE SECURITY LAYER   |
+-------------------------+
| Device Authentication:  |
| +-- Unique device       |
|     tokens             |
| +-- Certificate        |
|     validation         |
| +-- Token rotation      |
| +-- Device              |
|     fingerprinting     |
|                         |
| Communication Security: |
| +-- Encrypted channels  |
| +-- Message integrity   |
| +-- Replay attack       |
|     prevention         |
| +-- Command auth        |
+-------------------------+

+-------------------------+
| MONITORING & RESPONSE   |
+-------------------------+
| Security Monitoring:    |
| +-- Real-time event     |
|     detection          |
| +-- Behavioral anomaly  |
|     detection          |
| +-- Failed auth         |
|     tracking           |
|                         |
| Incident Response:      |
| +-- Automated alerts    |
| +-- Security incident   |
|     workflow           |
| +-- Forensic data       |
|     collection         |
|                         |
| Compliance & Audit:     |
| +-- Audit trail         |
|     maintenance        |
| +-- Security            |
|     assessments        |
| +-- Vulnerability       |
|     scanning           |
+-------------------------+
        \end{verbatim}
        \caption{Data Protection and Monitoring Systems}
        \label{fig:security-architecture-part2}
    \end{minipage}
\end{figure}

\section{Open Items and Future Work}
\label{sec:open-items-future-work}

Based on the current implementation and evaluation, we recommend the following enhancements:
\begin{itemize}
  \item \textbf{Database migrations}: Introduce Flyway/Liquibase for schema evolution across environments.
  \item \textbf{Indexes}: Add pragmatic indexes on high-traffic foreign keys and date/status columns (e.g., payments, service\_status, tamper/security logs).
  \item \textbf{Monetary consistency}: Standardize precision/scale and currency handling; consider enum for payment method.
  \item \textbf{Retention/archiving}: Define data retention and archiving for time-series tables (EnergyData, SecurityLog, TamperEvent).
  \item \textbf{Invariants}: Consider DB constraints to enforce business rules (e.g., a single active service status per installation if required).
  \item \textbf{Caching}: Add read-side caching for endpoints identified via profiling as hotspots.
  \item \textbf{E2E and load testing}: Add frontend e2e (Playwright) and basic load testing for ingestion and WebSocket broadcast.
\end{itemize}

\subsubsection{Security Implementation Details}
\label{subsubsec:security-implementation}

\paragraph{Authentication and Authorization Framework}
\begin{itemize}
    \item \textbf{JWT Token Structure}: RS256 algorithm with 24-hour expiration and refresh token rotation
    \item \textbf{Role-Based Access Control}: Hierarchical permission system with fine-grained control
    \item \textbf{Session Management}: Secure token storage with automatic cleanup and session timeout
    \item \textbf{Multi-Factor Authentication}: Framework ready for future MFA implementation
\end{itemize}

\paragraph{Data Protection and Encryption}
\begin{itemize}
    \item \textbf{At-Rest Encryption}: AES-256 encryption for sensitive database fields
    \item \textbf{In-Transit Encryption}: TLS 1.3 for all network communications
    \item \textbf{Key Management}: Secure key storage and rotation procedures
    \item \textbf{Password Security}: BCrypt with 12 salt rounds and complexity requirements
\end{itemize}

\paragraph{API Security Measures}
\begin{itemize}
    \item \textbf{Input Validation}: Comprehensive validation using Bean Validation annotations
    \item \textbf{Output Sanitization}: XSS prevention with proper encoding
    \item \textbf{CSRF Protection}: Synchronizer token pattern implementation
    \item \textbf{Rate Limiting}: Per-endpoint and per-user rate limiting policies
\end{itemize}

\section{Implementation}
\label{sec:implementation}

The backend is implemented in Java (Spring Boot), the frontend in Next.js/TypeScript, and the database in PostgreSQL. Key implementation decisions include modular service structure, REST API design, and JWT-based security. Data is validated and sanitized at all entry points.

Example: Energy reading processing (simplified):
\begin{verbatim}
public EnergyReading processReading(ReadingRequest req) {
    validateInstallation(req.getInstallationId());
    EnergyReading reading = createReading(req);
    save(reading);
    notifySubscribers(reading);
    return reading;
}
\end{verbatim}

\section{Testing}
\label{sec:testing-comprehensive}

Testing includes unit tests for business logic, integration tests for APIs, and end-to-end tests for user flows. Security and performance are validated with automated scripts. Example test:
\begin{verbatim}
@Test
public void processPayment_ValidRequest_Success() {
    PaymentRequest request = createValidRequest();
    PaymentResult result = service.processPayment(request);
    assertThat(result.isSuccess()).isTrue();
}
\end{verbatim}

\section{Deployment and Performance}

The system is deployed using Docker containers behind an NGINX proxy. The database is backed up regularly. Performance is optimized via query tuning, connection pooling, and asynchronous processing for high-volume data.

\subsubsection{Development Environment Setup}
\label{subsubsec:dev-environment-setup}

\paragraph{Prerequisites Installation}
\begin{verbatim}
# Java Development Kit 21
# Download from https://openjdk.org/projects/jdk/21/
# Verify installation
java --version
javac --version

# Node.js 20+ and npm
# Download from https://nodejs.org/
# Verify installation
node --version
npm --version

# Docker and Docker Compose
# Download from https://docker.com/
# Verify installation
docker --version
docker-compose --version

# Git version control
git --version

# Optional: PostgreSQL for local development
# Download from https://postgresql.org/
psql --version
\end{verbatim}

\paragraph{Quick Start with Docker (Recommended)}
\small
\begin{verbatim}
# 1. Clone the repository
git clone https://github.com/2001J/nebular-power.git
cd nebular-power

# 2. Create environment configuration
cp .env.template .env

# Edit .env file with your configuration:
# JWT_SECRET=your_32_character_secret_key_here
# JWT_EXPIRATION=86400000
# ADMIN_EMAIL=admin@yourdomain.com
# ADMIN_PASSWORD=your_secure_admin_password

# Generate JWT secret using OpenSSL
openssl rand -base64 32

# 3. Build and start all services
docker-compose up --build

# 4. Verify services are running
# Backend API: http://localhost:8080
# Frontend: http://localhost:3000
# Database: localhost:5432

# 5. Access the application
# Login with admin credentials from .env file
\end{verbatim}

\subsubsection{Production Deployment}
\label{subsubsec:production-deployment}

\paragraph{Production-Ready Docker Compose Configuration}
\small
\begin{verbatim}
version: '3.8'

services:
  postgres:
    image: postgres:15-alpine
    container_name: solar_postgres_prod
    environment:
      POSTGRES_DB: solar_energy_db
      POSTGRES_USER: ${DB_USER}
      POSTGRES_PASSWORD: ${DB_PASSWORD}
    volumes:
      - postgres_data_prod:/var/lib/postgresql/data
      - ./backups:/backups
    restart: unless-stopped
    networks:
      - solar_network
    ports:
      - "127.0.0.1:5432:5432"

  backend:
    build:
      context: .
      dockerfile: Dockerfile
      target: production
    container_name: solar_backend_prod
    depends_on:
      - postgres
    environment:
      SPRING_PROFILES_ACTIVE: production
      DB_HOST: postgres
      DB_NAME: solar_energy_db
      DB_USER: ${DB_USER}
      DB_PASSWORD: ${DB_PASSWORD}
      JWT_SECRET: ${JWT_SECRET}
      JWT_EXPIRATION: ${JWT_EXPIRATION}
      ADMIN_EMAIL: ${ADMIN_EMAIL}
      ADMIN_PASSWORD: ${ADMIN_PASSWORD}
    restart: unless-stopped
    networks:
      - solar_network
    healthcheck:
      test: ["CMD", "curl", "-f", "http://localhost:8080/actuator/health"]
      interval: 30s
      timeout: 10s
      retries: 3
      start_period: 60s

  frontend:
    build:
      context: ./solar_frontend
      dockerfile: Dockerfile
      target: production
      args:
        NEXT_PUBLIC_API_URL: ${NEXT_PUBLIC_API_URL}
    container_name: solar_frontend_prod
    depends_on:
      - backend
    restart: unless-stopped
    networks:
      - solar_network

  nginx:
    image: nginx:alpine
    container_name: solar_nginx_prod
    ports:
      - "80:80"
      - "443:443"
    volumes:
      - ./nginx/nginx.conf:/etc/nginx/nginx.conf:ro
      - ./ssl:/etc/nginx/ssl:ro
    depends_on:
      - backend
      - frontend
    restart: unless-stopped
    networks:
      - solar_network

volumes:
  postgres_data_prod:

networks:
  solar_network:
    driver: bridge
\end{verbatim}

\paragraph{Production Environment Variables}
\begin{verbatim}
# Database Configuration
DB_USER=solar_user
DB_PASSWORD=your_secure_database_password

# JWT Configuration
JWT_SECRET=your_32_character_jwt_secret_key_here
JWT_EXPIRATION=3600000  # 1 hour for production

# Admin Configuration
ADMIN_EMAIL=admin@yourdomain.com
ADMIN_PASSWORD=your_secure_admin_password

# API Configuration
NEXT_PUBLIC_API_URL=https://api.yourdomain.com

# Email Configuration (when implemented)
SMTP_HOST=smtp.yourdomain.com
SMTP_PORT=587
SMTP_USERNAME=noreply@yourdomain.com
SMTP_PASSWORD=your_smtp_password

# SSL Configuration
SSL_CERT_PATH=/etc/nginx/ssl/cert.pem
SSL_KEY_PATH=/etc/nginx/ssl/private.key

# Monitoring Configuration
LOGGING_LEVEL=INFO
METRICS_ENABLED=true
HEALTH_CHECK_ENABLED=true
\end{verbatim}

\paragraph{Database Backup and Recovery}
\begin{verbatim}
# Automated backup script
#!/bin/bash
# backup.sh

BACKUP_DIR="/backups"
DB_NAME="solar_energy_db"
DB_USER="solar_user"
TIMESTAMP=$(date +"%Y%m%d_%H%M%S")
BACKUP_FILE="$BACKUP_DIR/solar_backup_$TIMESTAMP.sql"

# Create backup
docker exec solar_postgres_prod pg_dump -U $DB_USER -d $DB_NAME > $BACKUP_FILE

# Compress backup
gzip $BACKUP_FILE

# Remove backups older than 30 days
find $BACKUP_DIR -name "solar_backup_*.sql.gz" -mtime +30 -delete

echo "Backup completed: ${BACKUP_FILE}.gz"

# Add to crontab for daily backups at 2 AM
# 0 2 * * * /path/to/backup.sh
\end{verbatim}

\subsubsection{Monitoring and Health Checks}
\label{subsubsec:monitoring-health}

\paragraph{Health Check Endpoints}
\begin{verbatim}
# Spring Boot Actuator endpoints
GET /actuator/health          # Overall application health
GET /actuator/health/db       # Database connectivity
GET /actuator/health/diskSpace  # Disk space availability
GET /actuator/metrics         # Application metrics
GET /actuator/prometheus      # Prometheus metrics format

# Custom health indicators
@Component
public class DatabaseHealthIndicator implements HealthIndicator {
    
    @Override
    public Health health() {
        try {
            // Check database connectivity and response time
            long startTime = System.currentTimeMillis();
            userRepository.count();
            long responseTime = System.currentTimeMillis() - startTime;
            
            if (responseTime > 1000) {
                return Health.down()
                    .withDetail("responseTime", responseTime + "ms")
                    .withDetail("status", "Slow database response")
                    .build();
            }
            
            return Health.up()
                .withDetail("responseTime", responseTime + "ms")
                .withDetail("status", "Database connection healthy")
                .build();
                
        } catch (Exception e) {
            return Health.down()
                .withDetail("error", e.getMessage())
                .build();
        }
    }
}
\end{verbatim}

\subsubsection{Debugging Techniques}
\label{subsec:debugging}

Effective debugging techniques include:

\begin{itemize}
    \item \textbf{Log Analysis}: Reviewing application logs for error messages and stack traces
    \item \textbf{Remote Debugging}: Attaching debuggers to running services
    \item \textbf{API Testing}: Using tools like Postman to test API endpoints in isolation
    \item \textbf{Database Queries}: Examining database queries and execution plans
    \item \textbf{Network Analysis}: Monitoring network traffic with tools like Wireshark
    \item \textbf{Profiling}: Using profilers to identify performance bottlenecks
\end{itemize}

Example of enabling remote debugging for backend services:

\begin{verbatim}
java -agentlib:jdwp=transport=dt_socket,server=y,suspend=n,address=5005 \
     -jar target/solar-0.0.1-SNAPSHOT.jar
\end{verbatim}

\subsection{Support Resources}
\label{subsec:support-resources}

The primary support resource for developers is the project's GitHub repository at \url{https://github.com/2001J/nebular-power}, which contains:

\begin{itemize}
    \item \textbf{Detailed README files}: Each module has a dedicated README file with specific implementation details
    \begin{itemize}
        \item ENERGY\_MONITORING\_README.md
        \item PAYMENT\_COMPLIANCE\_README.md
        \item SECURITY\_TAMPERING\_README.md
        \item SYSTEM\_MONITORING\_README.md
        \item USER\_MANAGEMENT\_README.md
    \end{itemize}
    \item \textbf{Technical Documentation}: Located in the /docs directory, organized by user and developer documentation
    \item \textbf{Source Code}: Well-commented code with clear structure following the organization described in this document
\end{itemize}

For specific issues or questions, developers should refer to the appropriate module README or examine the corresponding test files, which provide practical examples of component usage.

\section{Conclusion}
\label{sec:dev-conclusion}

The Solar Energy Monitoring and Control System represents a modern, scalable solution for monitoring and managing solar energy installations, with a particular focus on payment compliance and system security. The technical architecture and implementation details provided in this documentation should enable developers to understand, maintain, and extend the system effectively.

The modular monolithic architecture ensures that the system is maintainable and can evolve to meet changing requirements, while the comprehensive testing strategy helps maintain reliability and security. The documented APIs provide clear interfaces for integration with other systems, and the development guidelines ensure consistent, high-quality contributions.

As the system continues to evolve, the roadmap and architecture evolution plans provide a clear direction for future development, incorporating emerging technologies and addressing changing market needs. With the support resources and troubleshooting guidance provided, developers should be well-equipped to tackle any challenges that arise during development and operation.

By following the practices and principles outlined in this documentation, developers can contribute to the ongoing success and evolution of the Solar Energy Monitoring and Control System, helping to make solar energy more accessible and manageable for providers and customers alike.
